\documentclass[12pt]{article}
\usepackage[utf8]{inputenc}
\usepackage[T1]{fontenc}
\usepackage{hyperref}
\title{Kowaki Language Resources for K-12 Curriculum Development}
\author{Compiled by Claude Sonnet 5.5 (claude-sonnet-5-5) for Andr\'as Kornai}
\date{}
\begin{document}
\maketitle
\section{Introduction}
Kowaki (ISO 639-3 xow; Glottocode kowa1245) is a nearly extinct Papuan language of Madang Province, Papua New Guinea. It has about 25 cited speakers (2000) and is under pressure from the larger neighbouring Mauwake language. Glottolog lists its status as moribund. Almost nothing was found beyond catalogue entries.

\section{Primary Educational Resources}
None found.

\section{Linguistic and Reference Materials}
\begin{itemize}
\item Wikipedia, Kowaki language: \url{https://en.wikipedia.org/wiki/Kowaki_language}
\item Glottolog entry: \url{https://glottolog.org/resource/languoid/id/kowa1245}
\end{itemize}

\section{Religious and Community Resources}
The Joshua Project language page for xow gives only basic data: \url{https://joshuaproject.net/languages/xow}.

\section{Digital and Audio Resources}
None found. A search result mentioned a ``Kowai'' language on Umboi Island, but it was not confirmed to be the same language and is not used here.

\section{Curriculum Development Context}
No school programme, orthography or teaching materials were found. Any curriculum would have to begin with documentation work with the few remaining speakers.
\end{document}